% GENERATED FILE. Edit canonical lessons, then run npm run generate:books.
% canonical-source-hash: fnv1a64:7686a9ba811d62c6
% canonical-lessons: AR-C34-milh, AR-C34-lahm, AR-C34-jubn, AR-C34-taam

\chapter{Asking for Food}
\label{ch:ar-asking-for-food}
\hlchaptermodality{34}

\begin{chapteropening}
\textbf{By the end of this chapter:} \emph{I can ask for salt, meat, cheese and food in Arabic, and I can build the word for a restaurant myself out of a root and the ma- place shape.}

Name \ar{ط} as the emphatic partner of \ar{ت} beside \ar{ص} under \ar{س} and \ar{ض} under \ar{د}, then pour \ar{ط}-\ar{ع}-\ar{م} through the ma- place shape that already gave maktab and madhhab to generate maṭʿam, 'a restaurant', and ask for the chapter's foods with min faḍlik.
\end{chapteropening}

\section[milḥ]{\ar{ملح} (milḥ) — \textquotedblleft{}salt,\textquotedblright{} and the sailor hiding in it}

\label{lesson:AR-C34-milh}



\begin{quote}
A word for something you can hold in one hand — whose root stretches from a mineral to the open sea to a compliment.
\end{quote}

\subsection*{The letters in this word}
Nothing new, but one shape is worth a second look. \textbf{\ar{م}} (\emph{mīm}), \textbf{\ar{ل}} (\emph{lām}), \textbf{\ar{ح}} (\emph{ḥāʾ}) — all three join, so \textbf{\ar{ملح}} is one run of pen.

The \textbf{\ar{ح}} is the hook-and-tail letter that opened \textbf{\ar{حليب}} (\emph{ḥalīb}), but at the \textbf{end} of a word it grows its full tail and swings below the line: open at the start of \emph{ḥalīb}, closed and tailed at the end of \emph{milḥ}.

\textbf{\ar{قهوة}} made the same point at the other end of a word: \textbf{\ar{ة}} exists only in final position.

\begin{cousinweb}
\textbf{\ar{مِلْح}} (\emph{milḥ}) = \textquotedblleft{}\textbf{salt}.\textquotedblright{} Root \textbf{\ar{م}-\ar{ل}-\ar{ح}}, and this is a root that has gone travelling inside its own language:

\begin{itemize}
\raggedright
  \item \textbf{\ar{مَلّاح}} (\emph{mallāḥ}) — \textquotedblleft{}\textbf{a sailor}.\textquotedblright{} The man of the salt water. It wears the doubled \emph{lām} shape Arabic keeps for trades and habits.
  \item \textbf{\ar{مَليح}} (\emph{malīḥ}) — \textquotedblleft{}\textbf{handsome, pleasant}.\textquotedblright{} Salty, said of a person. Arabic reached for the same idea English reaches for when it calls a person \emph{the salt of the earth}.
\end{itemize}

Note \emph{malīḥ}'s shape: \textbf{\ar{فَعيل}} (\emph{faʿīl}), the same pattern that gave you \textbf{\ar{حَليب}} (\emph{ḥalīb}). A root plus that pattern hands back a quality.

Hebrew has \textbf{\he{מֶלַח}} (\emph{melaḥ}), \textquotedblleft{}salt,\textquotedblright{} unchanged — one more Semitic twin. English is \textbf{not} in this family: \emph{salt} comes down from Latin \textbf{sāl} and its Germanic cousins. And unlike \emph{sukkar}, nothing here was borrowed.
\end{cousinweb}

\subsection*{Your turn}
Say these aloud:

\begin{itemize}
\raggedright
  \item mīm, lām, ḥāʾ — \textquotedblleft{}milḥ,\textquotedblright{} salt
  \item ḥāʾ in two places — opening ḥalīb, closing milḥ with its tail
  \item the root's reach — mallāḥ, a sailor; malīḥ, handsome
  \item the faʿīl pattern twice — ḥalīb, malīḥ
  \item the Hebrew twin — melaḥ; and no English relative, only Latin sāl
  \item the borrowed one for contrast — sukkar, from Sanskrit by way of Persian
  \item the letters just behind — sīn, kāf, rā; qāf, hāʾ, wāw, tāʾ marbūṭa
\end{itemize}

\begin{tcolorbox}[breakable,colback=teal!4,colframe=teal!35!black,title={Before you move on}]
What else does \textbf{\ar{م}-\ar{ل}-\ar{ح}} give besides salt? (\emph{\textbf{Mallāḥ}}, a sailor, and \emph{\textbf{malīḥ}}, handsome.) Which pattern do \emph{malīḥ} and \emph{ḥalīb} share? (\emph{\textbf{Faʿīl}}.) Why does the \textbf{\ar{ح}} of \emph{milḥ} not look like the \textbf{\ar{ح}} of \emph{ḥalīb}? (\textbf{It is at the end of the word}, so it takes its full tailed shape.) Does English \emph{salt} belong to this family? (\textbf{No} — it descends from Latin \emph{sāl}.)
\end{tcolorbox}
\section[laḥm]{\ar{لحم} (laḥm) — \textquotedblleft{}meat,\textquotedblright{} and the town on the seam}

\label{lesson:AR-C34-lahm}



\begin{quote}
This word and its Hebrew twin are the same root — and they mean different foods. There is a town whose name proves it.
\end{quote}

\subsection*{The letters in this word}
Nothing new, and something pleasing: \textbf{\ar{لحم}} is \textbf{\ar{ملح}} read the other way about. The same shapes — \textbf{\ar{ل}} (\emph{lām}), \textbf{\ar{ح}} (\emph{ḥāʾ}), \textbf{\ar{م}} (\emph{mīm}) — in a different order, and so a different root and a different word.

That is the root system stated as plainly as it can be. The letters are not the word. The \textbf{order} is the word.

Here \emph{ḥāʾ} sits in the \textbf{middle}, so it wears its joined-both-sides shape: neither \emph{ḥalīb}'s open opening nor \emph{milḥ}'s tailed ending.

\begin{cousinweb}
\textbf{\ar{لَحْم}} (\emph{laḥm}) = \textquotedblleft{}\textbf{meat}.\textquotedblright{} Root \textbf{\ar{ل}-\ar{ح}-\ar{م}}.

Now the twin, and it is the best one in this book. Hebrew's \textbf{\he{לֶחֶם}} (\emph{leḥem}) is the same root — and it means \textbf{bread}.

Neither language changed the meaning. The old Semitic root named \textbf{the staple food}, and each people filled it with whatever theirs was: meat where herds mattered, bread where grain did.

English already owns the proof. The town is \textbf{\ar{بيت} \ar{لحم}} (\emph{Bayt Laḥm}) in Arabic and \textbf{\he{בֵּית} \he{לֶחֶם}} (\emph{Bēt Leḥem}) in Hebrew — one name, in which \emph{bayt} and \emph{bēt} both mean \textquotedblleft{}a house of.\textquotedblright{} English calls it \textbf{Bethlehem}: house of bread on one side of the seam, house of meat on the other.

And notice what Arabic did \textbf{not} do. It had this root available for bread and did not use it. Arabic's bread is \textbf{\ar{خبز}} (\emph{khubz}), built on \textbf{\ar{خ}-\ar{ب}-\ar{ز}}, \textquotedblleft{}to bake\textquotedblright{} — the baked thing, as \emph{ḥalīb} is the milked thing.
\end{cousinweb}

\subsection*{Your turn}
Say these aloud:

\begin{itemize}
\raggedright
  \item lām, ḥāʾ, mīm — \textquotedblleft{}laḥm,\textquotedblright{} meat
  \item the same shapes reordered — mīm, lām, ḥāʾ is \textquotedblleft{}milḥ,\textquotedblright{} salt
  \item ḥāʾ in the middle, joined on both sides
  \item the twin that shifted — laḥm is meat, Hebrew leḥem is bread
  \item the town on the seam — Bayt Laḥm, Bēt Leḥem, Bethlehem
  \item Arabic's own bread — khubz, from kh-b-z, to bake
  \item the salt root's reach — mallāḥ, malīḥ
  \item qahwa's letters and its feminine ending — qāf, hāʾ, wāw, tāʾ marbūṭa
\end{itemize}

\begin{tcolorbox}[breakable,colback=teal!4,colframe=teal!35!black,title={Before you move on}]
What separates \textbf{\ar{لحم}} from \textbf{\ar{ملح}}? (\textbf{Only the order of the letters} — the whole root system.) What does Hebrew \emph{leḥem} mean, and why is that no contradiction? (\textbf{Bread} — the root named the \textbf{staple food}.) Which English place name carries both? (\textbf{Bethlehem}.) What is Arabic's own word for bread, and where from? (\emph{\textbf{Khubz}}, from \textbf{\ar{خ}-\ar{ب}-\ar{ز}}, \textquotedblleft{}to bake.\textquotedblright{})
\end{tcolorbox}
\section[jubn]{\ar{جبن} (jubn) — \textquotedblleft{}cheese,\textquotedblright{} and where the root system stops}

\label{lesson:AR-C34-jubn}



\begin{quote}
The root system has been paying you back for chapters. This word is where it politely declines to.
\end{quote}

\subsection*{The letters in this word}
Nothing new. \textbf{\ar{ج}} (\emph{jīm}) is the third of the hook-and-tail family: the body of \textbf{\ar{ح}} (\emph{ḥāʾ}) and \textbf{\ar{خ}} (\emph{khāʾ}), marked by one dot \textbf{below}. Then \textbf{\ar{ب}} (\emph{bāʾ}), one dot below, and \textbf{\ar{ن}} (\emph{nūn}), one dot above.

Dots are the whole difference here, as they were between the \emph{ḥāʾ} of \textbf{\ar{لحم}} and the \emph{khāʾ} of \textbf{\ar{خبز}}.

\begin{cousinweb}
\textbf{\ar{جُبْن}} (\emph{jubn}) = \textquotedblleft{}\textbf{cheese}.\textquotedblright{} In everyday speech you will more often hear \textbf{\ar{جِبْنة}} (\emph{jibna}) — and that ending is the \textbf{\ar{ة}} you now recognise on sight, doing its usual work.

Now the honest part, because you have been taught to trust roots and this is where that trust needs a limit.

\textbf{\ar{جُبْن}} also means \textquotedblleft{}\textbf{cowardice},\textquotedblright{} and \textbf{\ar{جَبان}} (\emph{jabān}) is \textquotedblleft{}\textbf{a coward}.\textquotedblright{} Same letters, same order. So is a cheese a coward?

\textbf{No.} Arabic's dictionaries list these as \textbf{separate entries} — words spelled alike, not one idea branching. Compare a real root: \textbf{\ar{م}-\ar{ل}-\ar{ح}} carries salt into \emph{mallāḥ}, the sailor of salt water, and into \emph{malīḥ}, a salty compliment. You can feel the thread. Between cheese and cowardice there is none, and inventing one would be the mistake the milk story made of Aleppo.

Two words can share three letters by accident. Arabic has many roots and few letters.

English took none of this: \textbf{cheese} comes from Latin \textbf{cāseus}.
\end{cousinweb}

\subsection*{Your turn}
Say these aloud:

\begin{itemize}
\raggedright
  \item jīm, bāʾ, nūn — \textquotedblleft{}jubn,\textquotedblright{} cheese
  \item the hook family by dots — ḥāʾ none, khāʾ one above, jīm one below
  \item the everyday form — jibna, with the tied tāʾ
  \item the accident — jubn is also cowardice, jabān a coward
  \item the verdict — separate words, not one root branching
  \item a real branching — milḥ, mallāḥ, malīḥ
  \item meat and salt — laḥm, milḥ, the same shapes reordered
\end{itemize}

\begin{tcolorbox}[breakable,colback=teal!4,colframe=teal!35!black,title={Before you move on}]
What separates \textbf{\ar{ج}} from \textbf{\ar{ح}} and \textbf{\ar{خ}}? (\textbf{One dot below}.) What else do the letters of \emph{jubn} spell? (\textbf{Cowardice} — \emph{jabān} is a coward.) Are cheese and cowardice one root? (\textbf{No} — separate words that look alike.) Name a root that really does branch. (\textbf{\ar{م}-\ar{ل}-\ar{ح}} — \emph{milḥ}, \emph{mallāḥ}, \emph{malīḥ}.) What does the ending on \emph{jibna} tell you? (\textbf{Feminine}, as \emph{qahwa} is.)
\end{tcolorbox}
\section[ṭaʿām]{\ar{طعام} (ṭaʿām) — \textquotedblleft{}food,\textquotedblright{} and a restaurant you can build yourself}

\label{lesson:AR-C34-taam}



\begin{quote}
A new root, a new letter — and a word you will not have to be taught, because you already own the shape it comes in.
\end{quote}

\subsection*{The letters in this word}
One letter to name at last. \textbf{\ar{ط}} is \textbf{ṭāʾ}, the \textbf{emphatic} partner of \textbf{\ar{ت}} (\emph{tāʾ}): the same \emph{t}, said with the tongue flat, so the syllable goes dark around it.

Arabic keeps a small set of emphatic pairs, and you have met the others without being told they were a set:

\begin{itemize}
\raggedright
  \item \textbf{\ar{ص}} (\emph{ṣād}) is the emphatic \textbf{\ar{س}} (\emph{sīn}) — the \ar{ص} of \textbf{\ar{صباح}} (\emph{ṣabāḥ}).
  \item \textbf{\ar{ض}} (\emph{ḍād}) is the emphatic \textbf{\ar{د}} (\emph{dāl}) — the \ar{ض} of \textbf{\ar{من} \ar{فضلك}} (\emph{min faḍlik}).
  \item \textbf{\ar{ط}} (\emph{ṭāʾ}) is the emphatic \textbf{\ar{ت}} (\emph{tāʾ}) — this one.
\end{itemize}

Then \textbf{\ar{ع}} (\emph{ʿayn}), the throat sound already named, \textbf{\ar{ا}} (\emph{alif}) and \textbf{\ar{م}} (\emph{mīm}). \emph{Ṭāʾ} joins; \emph{alif} does not, so \textbf{\ar{طعام}} runs \textbf{\ar{طعا}}, then \textbf{\ar{م}}.

\begin{cousinweb}
\textbf{\ar{طَعام}} (\emph{ṭaʿām}) = \textquotedblleft{}\textbf{food}.\textquotedblright{} Root \textbf{\ar{ط}-\ar{ع}-\ar{م}}, whose verb \textbf{\ar{طَعِمَ}} (\emph{ṭaʿima}) means \textquotedblleft{}\textbf{he tasted},\textquotedblright{} and whose bare noun \textbf{\ar{طَعْم}} (\emph{ṭaʿm}) is \textquotedblleft{}\textbf{taste, flavour}.\textquotedblright{}

So Arabic has two ways to speak of eating. \textbf{\ar{أَكَلَ}} (\emph{akala}) is the plain act — he ate. \textbf{\ar{ط}-\ar{ع}-\ar{م}} goes through the tongue: tasting, and so the food that gets tasted.

Now take the \textbf{\ar{مَـ} place shape} you own. It made \textbf{\ar{مَكْتَب}} (\emph{maktab}), the place of writing, from \textbf{\ar{ك}-\ar{ت}-\ar{ب}}, and \textbf{\ar{مَذْهَب}} (\emph{madhhab}) from \textbf{\ar{ذ}-\ar{ه}-\ar{ب}}. Put \textbf{\ar{ط}-\ar{ع}-\ar{م}} through it:

\textbf{\ar{مَطْعَم}} (\emph{maṭʿam}) — the place of eating. \textbf{A restaurant.}

Nobody taught you that word. The pattern did.
\end{cousinweb}

\subsection*{Your turn}
Say these aloud:

\begin{itemize}
\raggedright
  \item ṭāʾ, ʿayn, alif, mīm — \textquotedblleft{}ṭaʿām,\textquotedblright{} food
  \item the emphatic set — ṣād under sīn, ḍād under dāl, ṭāʾ under tāʾ
  \item where you had them — ṣabāḥ, min faḍlik, now ṭaʿām
  \item the verb and the bare noun — ṭaʿima, he tasted; ṭaʿm, taste
  \item two ways to eat — akala the act, ṭaʿima the tasting
  \item the ma- place shape — madhhab, maktab, so maṭʿam, a restaurant
  \item the chapter's foods — milḥ, laḥm reversing its letters, jubn, ṭaʿām
  \item order them — \textquotedblleft{}laḥm, min faḍlik\textquotedblright{}; \textquotedblleft{}jubn, min faḍlik\textquotedblright{}
  \item and the odd old one — māʾ, water, whose plural miyāh keeps no pattern
\end{itemize}

\begin{tcolorbox}[breakable,colback=teal!4,colframe=teal!35!black,title={Before you move on}]
Which plain letter is \textbf{\ar{ط}} the emphatic partner of, and name the other pairs. (\textbf{\ar{ت}} — with \textbf{\ar{ص}} under \textbf{\ar{س}} and \textbf{\ar{ض}} under \textbf{\ar{د}}.) What does \textbf{\ar{ط}-\ar{ع}-\ar{م}} mean as a verb? (\textbf{To taste}.) Build \textquotedblleft{}restaurant\textquotedblright{} from it, and name the earlier words in that shape. (\emph{\textbf{Maṭʿam}} — like \emph{\textbf{maktab}} and \emph{\textbf{madhhab}}.) Ask for the chapter's foods politely. (\emph{\textbf{Milḥ}}, \emph{\textbf{laḥm}}, \emph{\textbf{jubn}}, \emph{\textbf{ṭaʿām}} — each with \emph{\textbf{min faḍlik}}, \textquotedblleft{}from your grace.\textquotedblright{}) Which everyday word keeps no pattern at all? (\emph{\textbf{Māʾ}} — its plural \emph{miyāh} is irregular.)
\end{tcolorbox}
